\documentclass[aspectratio=169]{beamer}
%\usetheme[height=8mm]{Rochester}
\usetheme{Boadilla}
\setbeamertemplate{blocks}[default]
\usepackage[french]{babel}
\usepackage[utf8]{inputenc}
\usepackage[T1]{fontenc}
\usepackage[sfdefault,light]{FiraSans}
%\renewcommand*\familydefault{\ttdefault} 
\usepackage{listings} 
\usepackage{color}
\usepackage{xcolor}
\lstdefinestyle{customc}{
	belowcaptionskip=1\baselineskip,
	breaklines=true,
	tabsize=2,
	frame=L,
	xleftmargin=\parindent,
	language=Python,
	showstringspaces=false,
	basicstyle=\footnotesize\ttfamily,
	backgroundcolor = \color{lightgray!15},
	keywordstyle=\bfseries\color{green!60!black},
	commentstyle=\itshape\color{purple!40!black},
	identifierstyle=\color{black},
	stringstyle=\color{orange},
	morekeywords={self}
}
\usepackage[style=numeric,
bibstyle=authoryear,
backend=bibtex]{biblatex}
\addbibresource{bibliography.bib}
\setbeamertemplate{navigation symbols}{}
\setbeamertemplate{footline}{\vspace*{1mm}\hfill
	\insertframenumber/\inserttotalframenumber\hfill\vspace*{1mm}}
\AtBeginSection[]{
	\begin{frame}[plain]
	\vfill
	\centering
	\begin{beamercolorbox}[sep=8pt,center]{title}
		\usebeamerfont{title}\insertsectionhead\par%
	\end{beamercolorbox}
	\vfill
\end{frame}
}

\begin{document}
\title{Most simple presentation}
\author{John Doe $^{\alpha}$\\~\\
\small{$^{\alpha}$ Some university}} 
\date{18 novembre 2019\\~\\\includegraphics[width=.5\textwidth]{images/ensta-couleur}} 
\frame[plain]{\titlepage} 

\begin{frame}{Frame title}
\begin{exampleblock}{Block title}
Hello world
\end{exampleblock}
\begin{block}{Block title}
Hello world \footfullcite{Feynman1941}
\end{block}
\begin{alertblock}{Block title}
Hello world \cite{Feynman1941}
\end{alertblock}
\end{frame}
\begin{frame}[fragile]{title}
\begin{lstlisting}[style=customc]
def __init__(self):
	self.name="toto"
def my_function(fname):
	print(fname + " Refsnes")
	return fname

my_function("Emil")
my_function("Tobias")
my_function("Linus") 
\end{lstlisting}
\end{frame}
\begin{frame}{Frame title}
	\begin{itemize}
		\item un
		\item deux
		\item trois
		\item quatre
	\end{itemize}
\end{frame}
\begin{frame}{Références}
\printbibliography
\end{frame}
\end{document}